% Figure: the second virial coefficient B(T) against temperature for the
% van der Waals form B(T) = b - a/(RT), showing the crossing from negative
% (attraction-dominated, low temperature) to positive (repulsion-dominated,
% high temperature) at the Boyle temperature T_B = a/(Rb). Reduced axes:
% B/b on the vertical, T/T_B on the horizontal, so the curve is universal.
\begin{figure}[htbp]
\centering
\begin{tikzpicture}
\begin{axis}[
  width=12cm, height=7.5cm,
  xlabel={reduced temperature $T/T_B$},
  ylabel={reduced second virial coefficient $B/b$},
  xmin=0.3, xmax=3.0, ymin=-2.2, ymax=1.1,
  axis lines=left,
  tick label style={font=\scriptsize},
  label style={font=\small},
  legend pos=south east,
  legend style={font=\scriptsize},
  clip=true,
]
% B/b = 1 - (T_B / T), since B = b - a/(RT) and T_B = a/(Rb)
\addplot[engblue, very thick, smooth, domain=0.35:3.0, samples=80] {1 - 1/x};
\addlegendentry{$B(T)=b-a/(RT)$}
% zero reference
\addplot[enggray, dashed, thick] coordinates {(0.3,0) (3.0,0)};
\addlegendentry{$B=0$ (ideal to first order)}
% Boyle temperature marker at T/T_B = 1, B = 0
\addplot[only marks, mark=*, engred, mark size=2.2pt] coordinates {(1,0)};
\node[anchor=north west, font=\scriptsize, engred] at (axis cs:1.03,-0.05)
  {Boyle temperature $T_B$};
% region labels
\node[font=\scriptsize\itshape, enggray] at (axis cs:0.6,-1.5)
  {attraction wins, $B<0$};
\node[font=\scriptsize\itshape, enggray] at (axis cs:2.2,0.75)
  {repulsion wins, $B>0$};
% high-T asymptote toward B = b
\addplot[engorange, dotted, thick] coordinates {(0.3,1) (3.0,1)};
\node[anchor=south east, font=\scriptsize, engorange] at (axis cs:3.0,1.0)
  {asymptote $B\to b$};
\end{axis}
\end{tikzpicture}
\caption[The second virial coefficient and the Boyle temperature]{The
second virial coefficient $B(T)$ in the van der Waals form $B=b-a/(RT)$,
plotted in reduced coordinates so the curve is the same for every substance.
At low temperature the molecules linger in one another's attractive wells,
the gas is easier to compress than ideal, and $B<0$. At high temperature the
molecules feel mostly the hard-core repulsion, the gas is harder to compress
than ideal, and $B>0$. The crossing at the Boyle temperature $T_B=a/(Rb)$ is
where the pairwise interactions cancel and a real gas obeys $pv=RT$ to first
order in density over a wide pressure range, behaving ideally not because it
is dilute but because attraction and repulsion balance. The van der Waals
value $T_B=(27/8)T_c$ overshoots the measured $T_B$ by the usual margin, but
the shape of the curve is correct.}
\label{fig:vol04ch02:boyle-virial-B}
\end{figure}
